% !TEX root = ../main.tex
\chapter{결론과 앞으로 할 일}
\label{ch:conclusion}

\ref{ch:discussion}장에서는 결과를 풀이하고, 이 연구가 목표를 이루었는지 판단했어요. \ref{ch:results}장의 표들이 그 기록이에요.

\dissertationSubheading[sec:summary-findings]{연구 질문에 대한 답}

이 논문은 엔도스랭크(EndorseRank)를 제안했어요. 엔도스랭크는 보증(믿고 맡김) 그래프(점과 화살표로 이은 그림) 위에서 계산하는 페이지랭크(PageRank)예요. 이 그래프는 마지막 허락(allowance, 토큰을 얼마까지 대신 써도 된다는 허락 가운데 가장 최근 것)으로 만들고, 화살표(엣지) 무게(가중치)는 적응형 가중 페이지랭크(AWP)의 것을 써요. 이 논문은 엔도스랭크를 AWP 자체\cend{3}, 두 종류의 화살표를 함께 걷는 결합 걷기, 그리고 이 걷기들이 매끄럽게 다듬는(평활하는) 원시 차수(화살표를 그냥 센 수)와 함께 평가했어요. 증거는 2025년 12월부터 2026년 5월까지 아비트럼 원(Arbitrum One)의 GMX~V2 지갑과 스펜더(spender, 허락을 받아 대신 쓰는 쪽)에서 나와요. 그리고 2026년 6월부터 8월까지의 자료로 등록 재현(계획을 먼저 공개해 두고 새 자료로 다시 해 보기)도 했어요.

연구 질문 RQ1은 엔도스랭크와 AWP가 한 표본의 지갑들에 순위를 다르게 매기는지 물었어요. 다르게 매겨요. 그래서 가설 H1은 뒷받침돼요(\ref{sec:rank-divergence}절). 차이의 많은 부분, 그리고 일치의 많은 부분이 엔도스랭크의 동점 덩어리(점수가 똑같은 지갑들의 무리)에서 나와요. 그래서 두 점수가 서로 다른 종류의 평판을 나타낸다는 주장은 주로 RQ4의 홀드아웃(나중 기간을 떼어 두고 맞혀 보는 시험)에 기대요.

RQ2는 같은 기간 안에서 각 점수가 허락 통계, 송금 통계와 얼마나 잘 맞는지 물었어요. 각 점수는 자기가 걷는 화살표를 따라가요(표~\ref{tab:alignment-hybrid}). 이것은 가설 H2를 좁은 뜻으로만 뒷받침해요. 대리 지표(대신 재는 값)가 그래프와 같은 이벤트(계약이 남기는 기록)에서 나오기 때문이에요. 또 엔도스랭크의 계수(일치 정도를 나타내는 수)는 동점이 정해요. 그래서 이 논문은 엔도스랭크의 허락 계수와 송금 계수 사이의 차이에서 아무 결론도 끌어내지 않아요.

RQ3은 각 점수를 계산하는 데 얼마나 드는지 물었어요. 비용은 화살표 수를 따라가요. 그래서 가장 듬성듬성한 그래프를 쓰는 엔도스랭크가 가장 싼 점수이고, 결합 걷기들은 AWP보다 더 많이 들어요(표~\ref{tab:benchmark-runtime}). 가설 H3은 뒷받침돼요. 확장 풀 단계(더 큰 주소 모음으로 넓힌 단계)에서는 화살표가 늘어날 때의 변화를 시험하지 못했어요. 두 번째 단계의 기록 꺼내기(확장 풀에 보탠 주소들의 기록 꺼내기)를 끝내지 못했기 때문이에요(표~\ref{tab:benchmark-scaling}).

RQ4는 $t_1$에 고정한 점수가 스펜더들에서 $t_2$의 새 허락(새로 생긴 주인--스펜더 허락)과 새 송금자(새로 송금을 보내온 주소)를 예측하는지(미리 맞히는지) 물었어요. 그리고 어떤 페이지랭크라도 자기가 다듬는 원시 차수보다 그것들을 더 잘 예측하는지 물었어요. 첫 번째는 맞아요. 그리고 점수가 어떤 라벨(나중에 채점에 쓰는 정답 값)을 미리 내다보는지는 화살표 종류가 정해요(표~\ref{tab:holdout-spenders}). 두 번째는 아니에요. 두 기간 모두에서, 원시 차수인 받은 허락 수(허락 입차수)는 새 허락을, 원시 차수인 송금 입차수(송금을 보내온 서로 다른 주소 수)는 새 송금자를 모든 걷기보다 더 잘 예측했어요(표~\ref{tab:holdout-diff}과~\ref{tab:fresh-posthoc}). 이 라벨들에서는 퍼뜨림(전파)이 그 자리에서 센 수보다 예측에 더 보태 주는 것이 없어요. 홀드아웃은 각 구성개념(점수가 재려고 하는 생각)의 시간에 따른 예측 타당도(나중 일을 미리 맞히는지)를 시험해요. 신용 위험 타당도는 시험하지 않은 채로 남겨 둬요.

RQ5는 두 종류의 화살표를 함께 걷는 결합 페이지랭크(C-PR, 두 그래프를 함께 걷는 점수)가 송금만 걷는 걷기보다 나은지 물었어요. 그리고 어느 한 층(화살표 한 종류로 만든 그래프)만 걷는 것보다 지갑 순위를 더 잘 매기는지 물었어요. 두 답 모두 “아니요”예요. 각각 미리 정한 규칙에 따른 답이에요. 9월 14일의 규칙에서, C-PR은 주된 부분의 두 조건을 모두 통과하지 못했고, 보조 부분의 허락 조건도 통과하지 못했어요(\ref{sec:coupled-results}절; 표~\ref{tab:holdout-diff}). 봄에 새 허락에서 송금 걷기보다 나았던 탐색적(미리 정하지 않고 살펴본) 결과는 등록한 6월부터 8월까지의 기간에서도 다시 나타났어요. 하지만 새 송금자에서 등록한 허용 폭(이만큼까지는 져도 된다고 정한 폭)보다 큰 손해가 있을 가능성을 지울 수 없었어요. 그래서 첫 번째 답도 “아니요”예요. AWP와 견주면 새 송금자에서의 손해가 분명했어요(표~\ref{tab:fresh-contrasts}). 씨앗을 쓰는 빼 보기 실험은 송금 걷기와 같은 결과를 냈어요. 허락 화살표는 걷기가 그 화살표를 따라 움직일 때 보탬이 되고, 씨앗으로 쓸 때는 보탬이 되지 않아요.

주장~C1--C5(\ref{sec:empirical-claims}절)는 이 결과들을 정식으로 적은 것이에요. 어떤 점수를 실제로 쓸지는 두 가지에 달려 있어요. 쓰는 사람의 질문이 허락(권한을 맡기는 일)에 관한 것인지 흐름(돈의 흐름)에 관한 것인지, 그리고 그 사람이 점수 두 개를 돌릴 수 있는지 하나에 매여 있는지예요.

\dissertationSubheading[sec:contribution-knowledge]{지식에 보탠 새로운 기여}

\ref{sec:contributions}절에서는 기여(이 연구가 새로 보탠 것) 네 가지를 들었어요. \ref{ch:literature}장에서 살펴본 문헌의 말로 하면, 이 논문은 다음을 밝혀요.

이 연구 전에는, 공개 장부(누구나 볼 수 있는 블록체인 기록) 위의 그래프 평판이 지갑에 순위를 매길 때 송금을 보거나\cend{3,4}, 거래 네트워크를 따라 위험이 퍼지는 모습을 보거나\cend{12}, 분명한 평가와 맡겨 둔 몫(스테이크)을 봤어요\cend{25,26}. 그리고 디파이(DeFi, 탈중앙 금융)에서 실제로 쓰는 평판 도구는 자격 증명(그 사람에 대해 확인해 준 사실)으로 주소 주인에게 점수를 매겼어요\cend{64}. 이 가운데 어느 것도 ERC-20(토큰 공통 규칙) 허락을 그래프의 화살표로 쓰지 않았어요. ERC-20 허락은 돈을 쓸 힘을 일부러 맡겼다는 온체인(블록체인 위의) 기록이에요\cend{18}. 또 어느 것도 온체인 평판 점수를 하나의 코호트(묶음)에서, 각 구성개념의 점검과 그 점수가 다듬는 원시 차수에 견주어 평가하지 않았어요. 신뢰구간(믿을 만한 범위), 차이 검사, 시간 홀드아웃을 갖춘 평가 말이에요. 이 논문은 허락 화살표와 그런 평가를 내놓아요. 페이지랭크와 AWP의 무게 매기기는 기존 문헌에서 가져왔어요\cend{3,15}. 엔도스랭크는 화살표만 바꾸고, 결합 연산자(두 층을 섞어 걷는 계산 규칙)는 걷기를 바꿔요.

새로운 지식은 세 부분으로 되어 있어요. 온체인 평판 점수가 무엇을 재는지는 그 점수가 쓰는 화살표에 달려 있어요. 그리고 이것은 하나의 코호트에서, 불확실성(얼마나 확실하지 않은지)을 수로 나타내면서 확인할 수 있어요. 허락 화살표로 만든 점수는 새로 고치는 비용이 더 싸고, 송금을 바탕으로 한 점수와 달라요. 이 점수는 두 기간 모두에서 스펜더들을 나중에 받는 보증에 따라 줄 세워요. 이 라벨에서 송금만으로 만든 점수들은 낮은 값에 머물러요. 다만 이 점수는 들어오는 흐름(유입)이나 트레이더(거래하는 사람)에 대해서는 정보를 거의 담고 있지 않아요. 그리고 문헌이 권한 두 가지 추가 단계는 이 자료에서 결과를 낫게 하지 못했어요. 퍼뜨림은 그래프가 앞으로 자라는 것을, 자기가 다듬는 원시 차수보다 더 잘 예측하지 못해요. 또 두 종류의 화살표를 무게를 달아 합친 것은 둘 중 더 나은 한 층을 이기지 못하고, 등록한 규칙에서는 송금만 걷는 걷기도 이기지 못해요. 이 주장들을 시험할 수 있게 해 주는 구성 타당도(점수가 재려는 것을 정말 재는지)의 틀\cend{22,33}도 기여의 한 부분이에요. 공개된 처리 과정(파이프라인)도 마찬가지예요. 이 처리 과정이 공개한 요약으로 다른 사람들이 모든 표를 확인할 수 있어요. 탈중앙 시스템이 신용 점수를 잘못 계산하는지는 모두가 관심을 가질 문제이기 때문이에요\cend{2}.

범위는 일부러 좁게 잡았어요. 이 논문은 온체인 평판이 담보를 대신할 수 있다는 것을 보여 주지 않아요\cend{1,9}. 결과가 아비트럼 원과 살펴본 기간 밖에서도 들어맞는다는 것도 보여 주지 않아요. 그리고 엔도스랭크가 시빌(가짜 주소를 많이 만들어 속이는 일)을 견딘다고 주장하지도 않아요. 모든 점으로 고르게 순간 이동하는(다시 출발하는) 방식에서는, 공격자 주소들로 만든 닫힌 농장(가짜 주소 무리)이 엔도스랭크를 싸게 부풀릴 수 있어요(\ref{sec:sybil-cost-model}절).

\dissertationSubheading[sec:future-work]{앞으로 할 일}

아래 제안은 저마다 \ref{ch:discussion}장에서 찾은 한계 하나에 맞춰져 있어요.

\begin{itemize}
\item 두 번째 단계의 기록 꺼내기 마치기. 지금까지 확장 주소 풀(더 큰 주소 모음)은 점(노드)의 개수를 바꾸는 데만 썼어요. 보탠 주소들의 허락 로그와 송금 로그(계약이 남기는 기록)를 꺼내면, 화살표 수에 따라 알고리즘의 실행 시간이 어떻게 늘어나는지 잴 수 있을 거예요. 그러면 RQ3에 대한, 범위를 좁힌 주장이 규모에 따른 변화에 대한 일반적인 주장이 될 거예요.
\item 다른 환경에서 결과 검증하기. 같은 분석을 다른 롤업(레이어 2 블록체인의 한 종류)과 레이어 1 네트워크(바탕이 되는 블록체인)에서 해 보면 도움이 될 거예요. 각 점수가 자기 구성개념과 맞는 것과 홀드아웃 결과가 아비트럼 원 밖에서도 이어지는지 보여 줄 수 있어요. 디파이는 종류가 매우 다양해서(이질적이라서)\cend{1,8}, 이것이 외적 타당도(바깥 자료로도 맞는지)를 위협하는 가장 큰 것이에요.
\item 기간을 늘리고 신용과 이어진 홀드아웃 쓰기. 관찰 기간이 더 길면 켄달의 타우($\tau$, 두 줄 세우기가 얼마나 같은 순서인지 나타내는 수) 계수가 시장 상황이 달라져도 안정적으로 유지되는지 알 수 있을 거예요. 주인(토큰 주인)들에 대한 신용과 이어진 홀드아웃은 대출 서비스가 코호트를 충분히 덮는 곳에서 할 수 있어요. 이 홀드아웃은 두 그래프 가운데 어느 하나라도 신용 위험 정보를 담고 있는지 보여 줄 거예요. 지금의 홀드아웃은 이것을 보여 주지 못해요.
\item 금액 무게. 원래 단위(소수점 조정을 하지 않은 가장 작은 단위)에 $b=1$을 쓰면, AWP의 금액 무게는 원래 단위로 몇 개 이상인 허락이나 송금을 모두 거의 한 번으로 쳐요. 가격으로 맞춘 금액(같은 돈 단위로 바꾼 금액)에서 $b$ 값을 정하면 크기가 다시 중요해질 수 있어요. 봄 홀드아웃은 그럴 수도 있음을 보여 줘요. 금액을 그대로 써서 무게를 단 허락 층이 두 스펜더 라벨 모두를 엔도스랭크보다 더 잘 줄 세웠거든요. 이런 선택은 그것으로 점수를 매기기 전에 미리 정해 두어야 해요.
\item 공격을 가정한 평가. 시빌 비용 모형은 계산으로 따진 분석이에요. 실제로 관찰한 엔도스랭크 그래프에서 이 논문의 페이지랭크 코드로 계산해 뒷받침했어요(\ref{sec:sybil-cost-model}절). 블록체인 위에서 실제로 공격을 해 보지는 않았어요. TrustRank처럼\cend{17} 순간 이동과 매달린 점(나가는 화살표가 없는 점)의 몫 다시 나누기를 미리 검사한 씨앗 집합(믿을 만하다고 고른 출발점들)으로만 제한하면, 닫힌 농장이 이용하는 고른 유입을 없앨 수 있을 거예요. 그리고 이 방법은 같은 공격 아래에서 평가해야 해요. 자전 허락(자기끼리 꾸며 낸 허락), 서로 짠 스펜더 농장, 계산된 허락 취소로 실제 레드팀 시험(일부러 공격해 보는 시험)을 하면, 실제 수수료 시장과 탐지 규칙 아래에서 공격이 비용에 비해 실제로 얼마나 효과적인지 잴 수 있을 거예요\cend{11}. 또 적극적인 공격자 앞에서도 시빌 보정(가짜 주소 문제를 고려해 고친) 대리 지표들이 쓸모 있는 정보를 주는지 시험할 수 있을 거예요.
\item 이동 단계에서 정해진 비율로 섞는 것을 넘어선 결합. 금액을 그대로 쓴 층에서, 모든 점으로 고르게 순간 이동하는 정해진 $\lambda$ 하나만 시도했어요. 점마다 다르거나 학습으로 정한 무게, 층마다 다른 감쇠 계수(화살표를 따라갈 확률), 두 정상 벡터(오래 걸었을 때 각 점에 머무는 비율)를 섞기, 엔도스랭크와 AWP처럼 무게를 단 층은 아직 시험하지 않았어요. 이런 설계는 무엇이든 점수를 계산하기 전에 규칙을 정해야 해요. 또 새 자료로 재현을 등록하고, 구성개념마다의 점검과 시간으로 나눈 증거(앞 기간으로 점수를 매기고 뒤 기간으로 확인한 증거)를 차수 기준선(비교 기준)과 견주어 보고해야 해요. 그리고 그래프 퍼뜨림을 속을 알 수 없는 학습 점수와 구별해 주는 성질, 곧 누구나 따져 볼 수 있다는 점을 지켜야 해요\cend{2,14}.
\item 차수 이기기. 자기 구성개념이 자라는 것을 맞히는 데서, $t_1$의 같은 화살표 종류의 원시 차수가 이 논문의 모든 페이지랭크보다 더 잘 예측했어요. 이 논문의 홀드아웃이 이 분야에 열어 둔 질문은 이것이에요. 새로 들어온 것을 센 수가 아닌 라벨에서, 퍼뜨린 점수가 그 자리에서 센 수보다 기간 밖에서 더 보탬이 될까요?
\item 더 풍부한 허락의 뜻. ERC-20 허락이 드문 곳에서는, 서비스마다 따로 있는 위임 로그를 화살표 정의에 보태면 더 많은 지갑을 덮을 수 있을지 몰라요. 이때도 같은 보고 원칙을 지켜야 해요.
\item 증명 도구와 함께 쓰기. 패스포트 방식의 관문(통과해야 하는 검사)\cend{64}은 한 주소가 다른 사람과 구별되는 한 사람의 것인지에 답해요. 엔도스랭크는 다른 사람들이 그 주소에 맡긴 권한으로 주소의 순서를 정해요. 순위를 매기기 전에 그런 관문으로 주소를 걸러 내는 것은 그럴듯한 쓰임새예요. 그리고 이것은 시빌인지 아닌지 라벨이 붙은 자료로, 관문 자체의 약점\cend{65,66}도 염두에 두고 평가해야 해요.
\item 실제 쓰임 연구. 엔도스랭크를 위험 오라클(블록체인 밖의 위험 정보를 들여오는 장치), 거버넌스(운영 규칙을 정하는 일) 도구와 합쳐 보면, 엔도스랭크의 실행 시간 특성이 실제 서비스에서 점수를 새로 고치는 빈도를 버틸 수 있는지 시험할 수 있어요. 수수료 등급이나 감시 빈도에 대한 통제된 실험(조건을 정해 두고 하는 실험)을 하면, \ref{sub:deploy-industry}절의 도입 가설을 증거로 바꿀 수 있어요. 담보 줄이기는 여전히 가설로 남아 있고, 이 연구는 그것을 주장하지 않아요.
\end{itemize}

\dissertationSubheading[sec:closing-reflection]{맺음말}

공격자가 있고 가명으로 움직이는 시장에서, 온체인 평판이 담보의 필요를 없애 주지는 않아요\cend{9}. 하지만 행동의 짜임새를 읽을 수 있게 해 줄 수는 있어요. 그리고 어떤 짜임새를 읽는지는 쓰는 화살표에 달려 있어요. 송금은 경제적인 흐름을 기록하고, 허락은 맡긴 권한을 기록해요. 페이지랭크는 어느 그래프에든 순위를 매길 거예요. 하지만 순위 하나가 두 구성개념 모두에 쓸모 있지는 않아요. 그리고 두 그래프를 함께 걷는 걷기는 어느 쪽에서도 알맞은 그래프 하나보다 낫지 않았어요. 주인들이 어떤 스펜더를 믿는지 빠르고 따져 볼 수 있는 신호가 필요한 서비스나 연구자라면, 여기서 살펴본 것과 같은 자료에서는 허락 화살표를 써야 해요. 받은 허락 수는 그 가장 싼 모양이에요. 엔도스랭크는 보증하는 쪽이 누구인지도 셈에 넣고 싶은 사람을 위한, 퍼뜨린 모양이에요. 다만 닫힌 농장이 엔도스랭크를 싸게 부풀릴 수 있다는 점을 염두에 두어야 해요(\ref{sec:sybil-cost-model}절). 송금 중심지(허브, 송금이 많이 모이는 곳)를 설명하고 그곳으로 들어올 흐름을 내다보려면 송금 화살표를 써야 해요. 이 그래프들만으로 거래 이익을 예측할 수 있는지는 이 자료로 결론을 낼 수 없는 질문이에요.
